\documentclass[11pt]{article}

\usepackage[margin=1in]{geometry}
\usepackage{amsmath}
\usepackage{amssymb}

\usepackage{graphicx}
\usepackage{bm}
\usepackage{booktabs}
\usepackage{array}
\usepackage{amsthm}
\usepackage{enumerate}
\usepackage[numbers,sort&compress]{natbib}
\usepackage{hyperref}
\usepackage{lineno}

\hypersetup{
  colorlinks=true,
  linkcolor=blue,
  citecolor=blue,
  urlcolor=blue
}

\newtheorem{theorem}{Theorem}
\newtheorem{proposition}{Proposition}
\newtheorem{corollary}{Corollary}
\newtheorem{lemma}{Lemma}

\begin{document}

\title{Conditional $\mathbb{Z}_n$ Amplitude Results for a Torus-Knot Parametrisation}

\author{Steve Bico Mujjabi, MD\thanks{msbico@gmail.com}\\
VuraLabs, Kampala, Uganda\\
Independent researcher\\
ORCID: \href{https://orcid.org/0009-0001-0556-5516}{0009-0001-0556-5516}\\
\url{https://github.com/bampita-bico/CDFD-Part-I-Release}}

\date{July 2026}

\maketitle

\begin{abstract}
This paper studies conditional $\mathbb{Z}_n$ amplitude results for a selected
torus-knot parametrisation. It does not establish a hierarchy of physical
particle families or a vacuum mechanism.

Two conditional algebraic results are proved. Given the stated
maximum-entropy modelling axiom, the Brannen amplitude coefficient is
$c_n=\sqrt{2}$. Given that coefficient and the selected parametrisation,
$\mathbb{Z}_3$ is the only family with a connected all-positive domain. These
results do not derive the axiom, select a physical knot, or assign particles.

The displayed family scales are numerical extrapolations that assume, without
derivation, that a Faddeev--Niemi scaling transfers to the CDFT functional.
They are not physical predictions. The higher-family parametrisations also lack
a fully positive amplitude domain.

Across the series, the public CDFD Runtime expresses this equilibrium vacuum
limit as a reusable flow-constraint state grammar. The calculations below use
that equilibrium limit only; adaptive departures are left for later dynamical
work rather than fitted in this manuscript. The paper-local supplementary
scripts remain the audited reproduction layer for the figures and tables in
this Part I archive.
\end{abstract}

\noindent\textbf{Keywords:} Constraint-Driven Field Theory, Constraint-Driven Flux Dynamics, adaptive operating ratio, vacuum structure, Mujjabi tests, reproducible physics, torus knots, particle families

\vspace{0.75em}
\linenumbers

\section*{Project Context}
This manuscript constitutes Part I (Paper 06) of the multi-part series
\textit{Constraint-Driven Flux Dynamics (CDFD)}. The underlying mathematical
architecture builds directly upon the concepts established in Paper 05 of this
series~\cite{MujjabiV2026}.

\section*{Symbol Conventions}

\noindent\textbf{CDFD physical-vacuum symbol convention.}
\begin{center}
{\small
\begin{tabular}{p{0.16\linewidth}p{0.33\linewidth}p{0.41\linewidth}}
\hline
Symbol & Part I meaning & Physical reading \\
\hline
$J$ or $\Phi$ & Driving flux or throughput & Transport current, field throughput, circulation drive, or generic flux; a subscript fixes the exact channel \\
$C$ & Active constraint or capacity burden & Vacuum transport capacity, geometric confinement, stiffness, boundary resistance, or load-bearing limit \\
$S$ & Surface responsiveness & Local ability of the vacuum or modeled surface to reconfigure under drive \\
$M_s$ & Structural memory & Retained geometry, phase history, stiffness history, or accumulated constraint state \\
$\Psi_s$ & Adaptive operating ratio $(J/C)\cdot S\cdot M_s$ or $(\Phi/C)\cdot S\cdot M_s$ & Local bookkeeping gauge for constrained operation; it is not a new universal constant \\
$R,a,\chi$ & Ring radius, core radius, and aspect ratio $R/a$ & Geometry used to connect vortex structure to dimensionless electromagnetic scales \\
$\rho_0,c,\kappa$ & Vacuum density scale, propagation speed, and transport stiffness & Medium parameters used in stability, pressure, and self-consistency calculations \\
\hline
\end{tabular}
}
\end{center}

\noindent\textbf{Greek-symbol convention.}
\begin{center}
{\small
\begin{tabular}{p{0.16\linewidth}p{0.33\linewidth}p{0.41\linewidth}}
\hline
Symbol & Local role & Interpretation rule \\
\hline
$\alpha$ & Fine-structure constant unless a local equation states otherwise & In the Part I physics papers, un-subscripted $\alpha$ means the electromagnetic fine-structure constant \\
$\beta,\gamma$ & Local response, stiffness, coupling, or correction coefficients & Equation-local coefficients; subscripts bind them to the named mode, channel, or approximation \\
$\theta,\phi,\Phi_c$ & Phase, orientation, or chirality angles & Used for torus-knot orientation, vacuum phase, or chiral phase geometry as defined in the nearby equation \\
$\delta,\epsilon$ & Perturbation, residual, tolerance, or small-offset scales & Local stability margins, numerical residuals, or experimental tolerances \\
$\lambda,\mu,\nu$ & Rate, mass, mode, or frequency labels & Meaning is fixed by the equation, subscript, and physics context where the symbol appears \\
$\rho,\sigma,\tau$ & Density, spread/noise, or time-scale terms & Local density, variance, stress, relaxation, or transport-memory quantities \\
$\Omega,\Lambda$ & State/domain space or integrated measure where defined & Reserved for explicitly defined state spaces, spectra, or synthesis-level quantities \\
\hline
\end{tabular}
}
\end{center}


\section{Introduction}

Papers~I--V of the CDFT series established the lepton sector:
\begin{enumerate}
  \item \textbf{Paper~I}~\cite{MujjabiI2026}: $\alpha = 1/137.036$ from vortex ring equilibrium.
  \item \textbf{Paper~II}~\cite{MujjabiII2026}: Koide $Q = 2/3$ as a $\mathbb{Z}_3$ theorem;
    leptons as trefoil $T(2,3)$ modes.
  \item \textbf{Paper~III}~\cite{MujjabiIII2026}: Topological selection of $T(2,3)$;
    chirality equilibrium $3\theta + \phi_c = \pi$.
  \item \textbf{Paper~IV}~\cite{MujjabiIV2026}: Master self-consistency equation for $\rho_0$;
    Faddeev--Niemi prediction $M_5 = 460.4$~MeV.
  \item \textbf{Paper~V}~\cite{MujjabiV2026}: $\phi_c = \pi - 2/3$ from $\mathbb{Z}_3$
    self-consistency; lepton masses from one input.
\end{enumerate}

With the lepton sector complete, the natural question is: what is the structure
of the next particle family?

Paper~IV predicted the energy scale $M_5 = 460.4$~MeV for the cinquefoil
$T(2,5)$ via Faddeev--Niemi scaling.
The present paper studies the general $\mathbb{Z}_n$ structure of the
$T(2,n)$ families, proves that $c_n = \sqrt{2}$ universally, identifies
why $\mathbb{Z}_3$ is the unique exactly solvable case, and states precisely
what is needed to determine the $T(2,5)$ mass spectrum.

\subsection{Hierarchy as capacity-shell quantization}

The hierarchy developed here is the Mujjabi hierarchy in a precise sense: each
$T(2,n)$ family is treated as a distinct capacity shell of the same constrained
transport medium. The universal coefficient $c_n=\sqrt{2}$ is not just an
algebraic convenience; it is the common amplitude scale required for the
families to share one operating grammar. The hierarchy claim can fail if the
higher shells do not map to external spectra or if each family requires its own
new fitted selector.


\section{The General $\mathbb{Z}_n$ Brannen Parametrisation}
\label{sec:brannen}

\subsection{Setup}

The $\mathbb{Z}_n$ analogue of the Brannen parametrisation~\cite{Brannen2006}
for the amplitudes of the $n$ modes of a $T(2,n)$ vortex is:
\begin{equation}
  A_k \;=\; 1 + c_n \cos\!\left(\theta + \frac{2\pi k}{n}\right),
  \quad k = 0, 1, \ldots, n-1\,,
  \label{eq:Zn_brannen}
\end{equation}
and the mode masses are $m_k = M A_k^2$.
The coefficient $c_n$ controls the relative amplitude of the $\mathbb{Z}_n$
oscillation.

\subsection{$\mathbb{Z}_n$ Fourier identities}

Two identities hold for the $\mathbb{Z}_n$ phase lattice $\{2\pi k/n\}$:
\begin{align}
  \sum_{k=0}^{n-1} \cos\!\left(\theta + \frac{2\pi k}{n}\right) &= 0
  \label{eq:ident_I} \\
  \sum_{k=0}^{n-1} \cos^2\!\left(\theta + \frac{2\pi k}{n}\right) &= \frac{n}{2}
  \label{eq:ident_II}
\end{align}
for any $\theta$ and any $n \geq 2$.
From~\eqref{eq:ident_I}: $\sum_k A_k = n$ (the sum of amplitudes is the
dimension $n$, independent of $\theta$).
From both identities: $\text{Var}(A_k) = c_n^2/2$.


\section{Theorem 1: The Universal Coefficient $c_n = \sqrt{2}$}
\label{sec:cn}

\begin{theorem}[Universal Brannen coefficient]
  \label{thm:cn}
  The Brannen amplitude coefficient for any $\mathbb{Z}_n$ family is
  $c_n = \sqrt{2}$, uniquely fixed by the maximum-entropy self-consistency
  condition $\mathrm{Var}(A_k) = \overline{A_k}^2$.
\end{theorem}

\begin{proof}
From identities~\eqref{eq:ident_I}--\eqref{eq:ident_II}:
\[
  \overline{A_k} = 1, \qquad \mathrm{Var}(A_k) = c_n^2/2.
\]
The maximum-entropy condition requires $\mathrm{Var}(A_k) = \overline{A_k}^2$:
\[
  c_n^2/2 = 1 \;\implies\; c_n = \sqrt{2}.
\]
Uniqueness: the condition is a quadratic in $c_n$ with unique positive root.
\end{proof}

\paragraph{Physical interpretation.}
The condition $\text{Var}(A_k) = \overline{A_k}^2$ is equivalent to:
the standard deviation of the amplitude distribution equals its mean.
The amplitude can reach zero (nodal line) or $2\overline{A_k} = 2$
(maximum) with equal probability under the $\mathbb{Z}_n$ phase average.
This is the condition of maximal relative uncertainty in the mode
amplitudes --- the vacuum selects the coefficient that neither suppresses
nor privileges any mode over the background.

\paragraph{Numerical verification.}
For any $n \geq 2$ and any test angle $\theta$, the ratio
$\text{Var}(A_k)/\overline{A_k}^2 = 1.000\ldots$ to machine precision
when $c_n = \sqrt{2}$ (Table~1 of the supplementary).


\section{The Generalised Koide Invariant and the Valid Domain}
\label{sec:koide_general}

\subsection{Derivation of $Q_n = 2/n$}

\begin{proposition}[General Koide invariant]
  \label{prop:Qn}
  For $m_k = M A_k^2$ with $A_k$ as in~\eqref{eq:Zn_brannen} and $c_n = \sqrt{2}$,
  the generalised Koide ratio satisfies $Q_n = 2/n$ provided all $A_k \geq 0$.
\end{proposition}

\begin{proof}
Using identity~\eqref{eq:ident_II} with $c_n = \sqrt{2}$:
\[
  \sum_k m_k = M \sum_k A_k^2 = M\!\left(n + c_n^2 \cdot \frac{n}{2}\right) = 2Mn.
\]
When all $A_k \geq 0$: $\sqrt{m_k} = \sqrt{M} A_k$, so
$\sum_k \sqrt{m_k} = \sqrt{M}\sum_k A_k = \sqrt{M}\cdot n$ (by identity~\eqref{eq:ident_I}).
Therefore:
\[
  Q_n \;\equiv\; \frac{\sum_k m_k}{\Bigl(\sum_k \sqrt{m_k}\Bigr)^2}
  \;=\; \frac{2Mn}{M n^2} \;=\; \frac{2}{n}\,.
\]
\end{proof}

\subsection{The valid domain}

The derivation requires all $A_k \geq 0$.
With $c_n = \sqrt{2}$, the amplitude $A_k = 1 + \sqrt{2}\cos(\cdot)$
reaches its minimum of $1 - \sqrt{2} \approx -0.414$ when the cosine
equals $-1$.
For $n = 3$: the minimum over all $k$ reaches $+0.293$ (positive) at
$\theta = 0$, and reaches zero at $\theta = \pm\pi/12 = \pm 15^\circ$.
So $\mathbb{Z}_3$ has a valid domain $\theta \in (-\pi/12, \pi/12)$.
The physical value $\theta = 2/9 \approx 12.73^\circ$ lies inside this domain.

For $n \geq 5$: no connected interval of $\theta$ exists where all $n$
amplitudes are simultaneously positive.
Numerically:
\begin{center}
\begin{tabular}{ccc}
\toprule
$n$ & $\max_\theta[\min_k A_k]$ & Valid domain? \\
\midrule
3  & $+0.293$ & Yes: $\theta \in (-\pi/12,\; \pi/12)$ \\
5  & $-0.144$ & No \\
7  & $-0.274$ & No \\
9  & $-0.329$ & No \\
11 & $-0.357$ & No \\
\bottomrule
\end{tabular}
\end{center}

\begin{figure}[ht]
  \centering
  \includegraphics[width=0.8\textwidth]{figures/fig6_0_panel.pdf}
  \caption{Minimum amplitude over all modes as a function of the phase $\theta$. For $n=3$, there is a domain where all amplitudes are positive. For $n \geq 5$, no such domain exists, obstructing the extension of the pure Koide formula.}
  \label{fig:valid_domain}
\end{figure}

\begin{theorem}[$\mathbb{Z}_3$ uniqueness]
  \label{thm:Z3_unique}
  $\mathbb{Z}_3$ is the unique torus knot family $T(2,n)$ for which:
  \begin{enumerate}[(a)]
    \item $c_n = \sqrt{2}$ (Theorem~\ref{thm:cn}),
    \item $Q_n = 2/n$ is constant on a connected valid domain
      (Proposition~\ref{prop:Qn}),
    \item the self-consistency condition $n\theta_n = Q_n$ fixes $\theta_n$
      within that domain (Paper~V).
  \end{enumerate}
  For $n \geq 5$, property~(b) fails: no connected valid domain exists.
\end{theorem}

\begin{proof}
  Property~(a) holds for all $n$ (Theorem~\ref{thm:cn}).
  Property~(b) requires a connected domain where all $A_k > 0$;
  the table above shows this fails for $n \geq 5$.
  Property~(c) requires (b) as a prerequisite: $n\theta = 2/n$ with
  $\theta \in (-\pi/12, \pi/12)$ gives $\theta_3 = 2/9$, which satisfies
  $\theta_3 < \pi/12$ ($2/9 \approx 0.222 < \pi/12 \approx 0.262$).
  No analogous closure exists for $n \geq 5$.
\end{proof}

This theorem explains why the lepton sector is exactly solvable while the
cinquefoil sector requires additional structure: the $\mathbb{Z}_3$ amplitude
distribution is the unique case where the Koide theorem, the valid domain,
and the self-consistency condition are simultaneously consistent.


\section{The Torus Knot Energy Hierarchy}
\label{sec:hierarchy}

\subsection{Topological selection rule}

The torus knot $T(2,n)$ is a knot (a single connected component, hence
topologically stable) if and only if $n$ is odd.
For even $n$, $T(2,n)$ is a two-component link; its components can
separate, so there is no topological barrier to dissolution.
Odd $n$ therefore labels the knot cases in this mathematical construction;
it does not establish physical particle families.

\subsection{Energy scaling from Faddeev--Niemi}

The Faddeev--Niemi result~\cite{FaddeevNiemi1997} applies to a specific
Skyrme-type functional. CDFT has not been shown to share that functional or to
identify its index with Hopf charge. The following is a conditional analogue:
\begin{equation}
  M_n \;=\; M_3 \left(\frac{n}{3}\right)^{3/4}\,,
  \label{eq:Mn}
\end{equation}
where $M_3 = 313.86$~MeV is the trefoil scale derived in Paper~V
($M_3 = m_e / A_{e,3}^2$ with $\theta_3 = 2/9$~rad).

\begin{table}[h]
  \caption{CDFT torus knot energy hierarchy.
    $M_3 = 313.86$~MeV from Paper~V.
    Conditional extrapolations from fitted $M_3$ and unverified Faddeev--Niemi transfer.}
  \label{tab:hierarchy}
  \begin{tabular}{ccccc}
    \toprule
    $n$ & Knot & $M_n$ (MeV) & Symmetry & Status \\
    \midrule
    3 & $T(2,3)$ trefoil     & 313.9 & $\mathbb{Z}_3$ & Papers I--V (complete) \\
    5 & $T(2,5)$ cinquefoil  & 460.4 & $\mathbb{Z}_5$ & This paper \\
    7 & $T(2,7)$ heptafoil   & 592.5 & $\mathbb{Z}_7$ & Future work \\
    9 & $T(2,9)$             & 715.4 & $\mathbb{Z}_9$ & Future work \\
    11& $T(2,11)$            & 831.6 & $\mathbb{Z}_{11}$ & Future work \\
    13& $T(2,13)$            & 942.7 & $\mathbb{Z}_{13}$ & Future work \\
    \bottomrule
  \end{tabular}
\end{table}

\begin{figure}[ht]
  \centering
  \includegraphics[width=0.8\textwidth]{figures/fig6_1_panel.pdf}
  \caption{Conditional numerical extrapolations under a Faddeev--Niemi scaling assumption.}
  \label{fig:energy_hierarchy}
\end{figure}

These are the energy \emph{scales} of each family.
The mass of the lightest mode of family $n$ is $M_n \cdot (A_{k,\min}^{(n)})^2$,
which requires knowing $\theta_n$ (open problem).


\label{sec:cinquefoil}

\subsection{What is established}

\paragraph{Energy scale.}
$M_5 = 460.4$~MeV is a conditional numerical extrapolation from fitted $M_3$.

\paragraph{Topology.}
$T(2,5)$ is a knot, hence topologically stable.
It supports $\mathbb{Z}_5$-symmetric modes (five lobes).

\paragraph{Amplitude coefficient.}
$c_5 = \sqrt{2}$ (Theorem~\ref{thm:cn}).

\paragraph{Hadron scale context.}
$M_5 = 460.4$~MeV lies between the kaon ($\sim 494$~MeV) and the
eta meson ($\sim 548$~MeV).
We do not claim the cinquefoil corresponds to a specific known particle:
the mass spectrum depends on $\theta_5$ (open problem~1 below).

\subsection{The structural obstruction for $n \geq 5$}

With $c_5 = \sqrt{2}$, the five amplitudes
$A_k = 1 + \sqrt{2}\cos(\theta + 2\pi k/5)$ for $k = 0,\ldots,4$
are never all positive simultaneously (see the valid-domain table in
Section~\ref{sec:koide_general}).
The minimum amplitude over all five is always negative,
with maximum value $-0.144$ at $\theta = 0$.

This means:
\begin{enumerate}
  \item $Q_5 = 2/5$ does not hold as a universal constant
    (Proposition~\ref{prop:Qn} fails).
  \item The step $\sum_k \sqrt{m_k} = \sqrt{M}\sum_k A_k$ is invalid
    because some $A_k < 0$ requires $\sqrt{m_k} = \sqrt{M}|A_k|$.
  \item The self-consistency condition $5\theta_5 = Q_5 = 2/5$
    cannot be applied without first establishing the correct $Q_5$.
\end{enumerate}

\subsection{The modified Koide ratio for $n \geq 5$}

When some $A_k < 0$, define $\tilde{Q}_n(\theta)$ using $|A_k|$:
\begin{equation}
  \tilde{Q}_n(\theta) \;\equiv\;
  \frac{\sum_k A_k^2}{\Bigl(\sum_k |A_k|\Bigr)^2}\,.
  \label{eq:Q_tilde}
\end{equation}
For $\mathbb{Z}_3$ in its valid domain: $\tilde{Q}_3 = Q_3 = 2/3$ (constant,
since all $A_k > 0$ there).
For $\mathbb{Z}_5$: $\tilde{Q}_5(\theta)$ varies from $0.294$ to $0.326$,
far from the putative value $2/5 = 0.400$.

\begin{figure}[ht]
  \centering
  \includegraphics[width=0.8\textwidth]{figures/fig6_2_panel.pdf}
  \caption{The modified ratio $\tilde Q_n(\theta)$ for $n=3$ and $n=5$. While $\tilde Q_3$ is a strict constant ($2/3$), $\tilde Q_5$ varies and does not meet the $2/5$ target.}
  \label{fig:q_tilde}
\end{figure}

The function $\tilde{Q}_5(\theta)$ is not constant.
This means: no analogue of the Koide self-consistency argument from Paper~V
applies to $\mathbb{Z}_5$ without additional structure.
A derivation of $\theta_5$ requires either:
(a) a physical principle beyond the $\mathbb{Z}_n$ algebraic structure alone, or
(b) a modified Brannen parametrisation for the $T(2,5)$ family that avoids
    negative amplitudes.


\section{The CDFT Torus Knot Picture}
\label{sec:picture}

The torus knot hierarchy $T(2,n)$ for odd $n$ provides a candidate ordered
sequence of CDFT particle families.
Table~\ref{tab:hierarchy} lists the energy scales.

The key structural divide is at $n = 3$:
\begin{itemize}
  \item $n = 3$ (trefoil, leptons): $c_3 = \sqrt{2}$, valid domain exists,
    $Q_3 = 2/3$ constant, self-consistency closes the internal system.
  \item $n \geq 5$: $c_n = \sqrt{2}$, \emph{no} valid domain, $Q_n \neq 2/n$
    (as a constant). Requires new approach for mass spectrum.
\end{itemize}

The energy scale hierarchy is fixed once the model assumptions are accepted.
The mass spectra of $n \geq 5$ families are open.

\begin{table}[h]
  \caption{Torus knot $T(2,n)$ topological classification.}
  \label{tab:classification}
  \begin{tabular}{cccc}
    \toprule
    $n$ & Topological type & Persistent? & CDFT family \\
    \midrule
    1 & Unknot & No & (none) \\
    2 & Hopf link & No & (none) \\
    3 & Trefoil knot & \textbf{Yes} & Leptons (Papers I--V) \\
    4 & Solomon link & No & (none) \\
    5 & Cinquefoil knot & \textbf{Yes} & This paper \\
    7 & Heptafoil knot & \textbf{Yes} & Future \\
    $2k$ & $k$-component link & No & (none) \\
    $2k+1$ & Torus knot & \textbf{Yes} & Predicted \\
    \bottomrule
  \end{tabular}
\end{table}


\section{Open Problems}
\label{sec:open}

\begin{enumerate}

\item \textbf{The $T(2,5)$ chirality phase $\theta_5$.}
  The cinquefoil has $\mathbb{Z}_5$ symmetry and five modes.
  Determining $\theta_5$ requires either:
  (a) a physical principle that selects a preferred orientation of the
  $\mathbb{Z}_5$ oscillation in the vacuum background, or
  (b) a modified amplitude form for $T(2,5)$ that ensures all $A_k \geq 0$.

\item \textbf{The $T(2,5)$ mass spectrum.}
  With $\theta_5$ known, the five mode masses are
  $m_k^{(5)} = M_5 |A_k^{(5)}|^2$.
  Identification of these five modes with known particles
  (at the $\sim 460$~MeV scale: kaons, etas, or new states)
  depends on whether the cinquefoil modes are confined (quark-like)
  or free (lepton-like).

\item \textbf{Why is $\mathbb{Z}_3$ exactly solvable?}
  Theorem~\ref{thm:Z3_unique} identifies the structural reason
  ($n = 3$ uniquely has a valid domain with $c_n = \sqrt{2}$),
  but a deeper question remains: why does the CDFT vacuum select
  a ground state where the first stable knotted excitation ($T(2,3)$)
  is exactly the one where the Koide self-consistency is possible?

\item \textbf{The general $n$-family chirality equilibrium.}
  Paper~III derived $3\theta + \phi_c = \pi$ from the $\mathbb{Z}_3$
  chirality interaction $\lambda_c \cos(3\theta + \phi_c)$.
  The $\mathbb{Z}_5$ analogue would be $5\theta_5 + \phi_{c,5} = \pi$
  (or some other value).
  Whether this generalises directly and what determines $\phi_{c,5}$
  is an open problem.

\end{enumerate}


\section{Scope Boundary and Conclusion}

Under the stated amplitude ansatz and maximum-entropy modelling axiom, the
coefficient and $\mathbb{Z}_3$ domain results are conditional mathematical
results. They do not derive the axiom, a vortex energy functional, or a physical
particle interpretation. The family-scale table is an unvalidated extrapolation.
This paper must not be cited as a universal hierarchy, an explanation of 137,
or a particle-spectrum prediction.


\section*{Supplementary Material}

This manuscript is accompanied by release-archive supplementary material in
the canonical Part I tree. The relevant paper-local Python scripts, numerical checks, generated figure outputs,
tabulated data, figure assets, and environment notes are maintained under
\texttt{notebooks/}, \texttt{outputs/}, \texttt{figures/}, and
\texttt{REPRODUCIBILITY.md} in \texttt{Part\_I\_Fundamental\_Physics/}.


\section*{Data and Code Availability}

The release-archive supplementary material is included in the archived Part I
release on Zenodo at \url{https://doi.org/10.5281/zenodo.20250820}. The current version DOI is assigned by Zenodo after each GitHub release. Zenodo is the
permanent release archive, and the public source archive is available on GitHub at
\url{https://github.com/bampita-bico/CDFD-Part-I-Release}. The broader public CDFD
Runtime is separate reusable software infrastructure; the numerical values reported
here are reproduced from this paper's Supplementary Material.


\section*{Acknowledgments}
The author used AI-assisted tools for formatting checks and code-verification
support. The scientific claims, interpretations, and final manuscript decisions
are the author's responsibility.

The author is indebted to the foundational work of
L.~Faddeev and A.J.~Niemi~\cite{FaddeevNiemi1997} (knots and particles),
R.A.~Battye and P.M.~Sutcliffe~\cite{BattyeSutcliffe1998}
(knot soliton energy spectra),
D.~Rolfsen~\cite{Rolfsen1976} (torus knot classification),
and Y.~Koide~\cite{Koide1983} (lepton mass formula)
--- whose work this paper generalises to arbitrary $T(2,n)$ families.

\section*{Declarations}

\subsection*{Funding}
This research received no external institutional funding.

\subsection*{Author Contributions}
Steve Bico Mujjabi: conceptualization, methodology, formal analysis, software,
validation, visualization, writing--original draft, and writing--review and
editing.

\subsection*{Conflicts of Interest}
The author declares no conflicts of interest.

\section*{References}
\bibliographystyle{plainnat}
\bibliography{references}

\end{document}
